\section{Workload Definition}

Workloads are defined precisely to ensure experimental reproducibility and to simulate
realistic pipeline execution resource consumption patterns. Six concurrency levels are
tested to map the full degradation curve rather than sampling only the endpoints.

\begin{table}[H]
\centering
\caption{Defined Workload Levels}
\label{tab:workloads}
\begin{tabularx}{\textwidth}{lXXXX}
\toprule
\textbf{Level} & \textbf{Concurrent Jobs} & \textbf{Per-Job CPU} & \textbf{Per-Job Memory} & \textbf{Workload Type} \\
\midrule
L1 (Baseline)   &  1 & 0.25 CPU & 256 MB & Single job, no contention \\
L2 (Low)        &  2 & 0.25 CPU & 256 MB & Light concurrent load \\
L3 (Medium)     &  4 & 0.25 CPU & 256 MB & Moderate concurrent load \\
L4 (High)       &  8 & 0.25 CPU & 256 MB & Heavy concurrent load \\
L5 (Stress)     & 12 & 0.25 CPU & 256 MB & Near-capacity load \\
L6 (Overload)   & 16 & 0.25 CPU & 256 MB & Expected failure zone \\
\bottomrule
\end{tabularx}
\end{table}

Each workload level is run \textbf{three times} per environment. Results are reported as
mean and standard deviation across runs to account for natural system variability. The
synthetic pipeline job performs a fixed-duration computation (matrix operations via NumPy)
that produces deterministic resource consumption, making results reproducible across
environments.
